% --- Templates -------------------------------------------------------------
%\includegraphics{}

% Dokumentklasse --------------
\documentclass[12pt,naustrian,a4widepaper]{scrartcl}   
% article style
%   - 11pt Schriftgroesse
%   - new austrian (neue Rechtschreibung)
%   - Papierformat A4
%   - pdf-hyperlinks

% Packages --------
%\RequirePackage{hgblistings}
\usepackage[utf8]{inputenc}  % fuer Umlaute, äöü
\usepackage[T1]{fontenc}
\usepackage{a4wide}
\usepackage{times}      % Times Schriften (zusammen mit fontencoding, s.o.)
\usepackage{babel}
\usepackage{graphicx}	  % für das Einbinden von Grafiken
\usepackage{color}      % für färbigen Text
\usepackage{framed}     % für (Text-) Rahmen
\usepackage{listings}   % für den Sourcecode
\usepackage{fancyhdr}   % für Kopf- und Fusszeilen
\usepackage{pdfpages}	% pdf import fuer angabe

\pagestyle{fancy}       % Kopf- / Fusszeile aktivieren
\lstset{
  language={C++},
  basicstyle=\footnotesize\ttfamily,
  keywordstyle=\color{blue},%\bfseries,
  commentstyle=\color{green},
  frame=single,
  xleftmargin=0pt,
  framexleftmargin=0pt,
  framexrightmargin=0pt,
  linewidth=\dimexpr\textwidth+3cm\relax,
  breaklines=true,
  tabsize=3,
  numbers=left, numberstyle=\tiny, stepnumber=1, numbersep=5pt
}

% Seitenspiegel -------------
\typearea{8}	% Festlegung des Seitenspiegels gem. Koma. 4..groß, 9..klein

% Kopfzeile ---------
\lhead{{\footnotesize{Boris Dobrijevic, Maximilian Ebert}}}   %  (links)
\rhead{{\footnotesize{Seite \thepage}}}      %  (rechts)

% Fusszeile ---------
\lfoot{}  % links
\cfoot{}  % mitte 
\rfoot{}  % rechts

% Absatzformatierung ------------------
\setlength{\parindent}{0cm}   % Einrückung der 1. Zeile jedes Absatzes
\setlength{\parskip}{10pt}    % Abstand zwischen den Absätzen

% Package für Hyperlinks (mit pdf-Optionen) -----------------------------------------
\usepackage[
urlcolor=blue,		% blaue weblinks
linkcolor=black,	% interne Links sind schwarz
colorlinks=true,        % links werden eingefürbt
pdfstartview=FitH,      % PDF-Anzeige: Fensterbreite
pdfborder={0 0 0},      % keine Umrandung um links
pdftitle   ={Systemdokumentation},% Referenzen in der Pdf-Datei
pdfauthor  ={Boris Dobrijevic, Maximilian Ebert},
pdfsubject ={Systemdoku},
pdfcreator ={Der Creator},
pdfproducer={Der Producer},
pdfkeywords={Dokumentation, Systemdokumentation}
]{hyperref}

% =================================================================================================
% =================================================================================================
% Beginn des Dokumentes
\begin{document}
\selectlanguage{naustrian}   % oder "american" für engl. Texte

% Titelblatt ----------
\title {\vspace{1cm}
       \includegraphics[width=8cm]{./Images/FhOOeLogoOkt2009_HSD_Rot_pastell}\\
       \vspace{2cm}
       {\textbf{Systemdokumentation\\Übung 02}}\\
       \vspace{5mm}
       {\small{Version 1.0}}\\
       \vspace{5mm}
}

\author{\small{Boris Dobrijevic, Ebert Maximilian}}
\date  {\small{Hagenberg, \today}}
%\maketitle

\includepdf[pages={1-2}]{Uebung_4_RTO.pdf}
\clearpage

% Inhalts-, Tabellen- und Bildverzeichnis (werden generiert) ----------------------------------------------------------
\tableofcontents

\clearpage
% =================================================================================================
% Kapitel 1 ---------------------
\section{Aufgabe A: Erweiterung Mutex}

% =================================================================================================
% Kapitel 2
\section{Aufgabe B: Refaktoring und Verbesserung des Quellcodes}


% #################################################################################################
%      Quellcode
% #################################################################################################
\clearpage
\section{Quellcode}
% \subsection{main.c}
% \lstinputlisting{../../Uebung3/src/main.c}



\end{document}
